\cleardoublepage
\chapter{KORPUS PEDOMAN KLINIS}
\label{lamp:korpus}

Lampiran ini mendaftar dokumen yang menyusun basis pengetahuan sistem sebagaimana pembentukannya diuraikan pada Subbab~\ref{subsec:bab5-pembentukan-basis-pengetahuan}.
Daftar ini disertakan agar dua hal dapat diperiksa pembaca. Pertama, seluruh potongan yang dapat dikembalikan sistem berasal dari terbitan resmi lembaga profesi,
bukan dari ringkasan yang disusun sendiri. Kedua, tiap sitasi pada keluaran sistem dapat ditelusuri sampai ke halaman cetak dokumen sumbernya, yaitu kebutuhan
KNF-08 pada Subbab~\ref{sec:bab3-analisis-kebutuhan}. Daftar dokumennya ditunjukkan pada Tabel~\ref{tab:tabel_D1_korpus}.

\begingroup

\setlength{\LTcapwidth}{\textwidth}

\small

\captionsetup{font=normalsize}

\setlength{\tabcolsep}{4pt}

\begin{longtable}{@{}>{\raggedright\arraybackslash}p{1.000cm}>{\raggedright\arraybackslash}p{2.350cm}>{\raggedleft\arraybackslash}p{1.000cm}>{\raggedright\arraybackslash}p{7.075cm}>{\raggedleft\arraybackslash}p{1.450cm}@{}}

\caption{Dokumen pedoman klinis penyusun basis pengetahuan.}

\label{tab:tabel_D1_korpus} \\

\toprule

\textbf{Kode} & \textbf{Lembaga} & \textbf{Tahun} & \textbf{Judul dokumen} & \textbf{Halaman} \\

\midrule

\endfirsthead



\multicolumn{5}{@{}l}{\small\emph{Tabel \thetable{} (lanjutan)}} \\

\toprule

\textbf{Kode} & \textbf{Lembaga} & \textbf{Tahun} & \textbf{Judul dokumen} & \textbf{Halaman} \\

\midrule

\endhead



\bottomrule

\endlastfoot



KB-01 & IDAI UKK Endokrinologi Anak dan Remaja & 2015 & Konsensus Nasional Pengelolaan Diabetes Melitus Tipe-1 Edisi Ketiga (revisi 2021) & 129 \\



KB-02 & PERKENI & 2021 & Pedoman Pemantauan Glukosa Darah Mandiri 2021 & 49 \\



KB-03 & PERKENI & 2021 & Pedoman Petunjuk Praktis Terapi Insulin pada Pasien Diabetes Melitus 2021 & 70 \\



KB-04 & PERKENI & 2023 & Tatalaksana Pasien dengan Hiperglikemia di Rumah Sakit & 105 \\



KB-05 & ADA dan EASD & 2021 & The Management of Type 1 Diabetes in Adults: A Consensus Report. Diabetologia 64(12):2609-2652 & 44 \\



KB-06 & International Consensus (ATTD) & 2019 & Clinical Targets for Continuous Glucose Monitoring Data Interpretation. Diabetes Care 42:1593-1603 & 11 \\



KB-07 & ISPAD & 2022 & ISPAD CPCG 2022 Chapter 10: Nutritional Management & 25 \\



KB-08 & ISPAD & 2022 & ISPAD CPCG 2022 Chapter 11: Diabetic Ketoacidosis and Hyperglycemic Hyperosmolar State & 22 \\



KB-09 & ISPAD & 2022 & ISPAD CPCG 2022 Chapter 12: Assessment and Management of Hypoglycemia & 19 \\



KB-10 & ISPAD & 2022 & ISPAD CPCG 2022 Chapter 13: Sick Day Management & 14 \\



KB-11 & ISPAD & 2022 & ISPAD CPCG 2022 Chapter 14: Exercise in Children and Adolescents with Diabetes & 32 \\



KB-12 & ISPAD & 2022 & ISPAD CPCG 2022 Chapter 25: Management in Limited Resource Settings & 23 \\

\end{longtable}



\noindent\tabnotestyle Seluruhnya berjumlah 12 dokumen dengan total 543 halaman cetak. Kolom halaman menyatakan jumlah halaman dokumen sumber, bukan jumlah potongan hasil pemenggalan.\normalsize

\endgroup



Korpus terdiri atas empat kelompok penerbit. Konsensus IDAI dan tiga pedoman PERKENI mewakili tata laksana pada konteks Indonesia, termasuk pemantauan glukosa
mandiri, modalitas sasaran penelitian ini. Pernyataan bersama ADA dan EASD serta konsensus ATTD mengenai target glikemik melengkapi rujukan internasional.
Enam bab pedoman ISPAD mencakup keadaan yang paling sering muncul pada penilaian sistem, yaitu nutrisi, ketoasidosis,
hipoglikemia, manajemen hari sakit, aktivitas fisik, dan pengelolaan pada lingkungan bersumber daya terbatas.

Pemilihan mengikuti kriteria pada Subbab~\ref{subsec:bab4-kurasi-basis-pengetahuan}, yaitu terbitan resmi, penomoran halaman yang tetap, dan relevan dengan
keadaan glukosa yang dikenali sistem. Dokumen yang tidak memenuhi ketiganya tidak dimasukkan meskipun topiknya berdekatan.

Setiap dokumen disimpan bersama medan penggeser halaman pada manifes korpus. Medan tersebut diperlukan karena pada sebagian dokumen nomor halaman cetak berbeda
dari nomor halaman berkas digitalnya, sedangkan sitasi yang ditampilkan sistem harus menunjuk halaman cetak agar dapat diperiksa pembaca pada dokumen aslinya.
Jumlah potongan hasil pemenggalannya dilaporkan pada Subbab~\ref{subsec:bab5-pembentukan-basis-pengetahuan}.
